% !TEX program = xelatex
% À compiler sur Overleaf avec le compilateur XeLaTeX
% (Menu -> Compiler -> XeLaTeX) pour avoir la police Verdana.

\documentclass[12pt,a4paper]{article}

\usepackage{fontspec}
% Verdana si disponible, sinon une police sans-serif proche
\IfFontExistsTF{Verdana}{\setmainfont{Verdana}}{\setmainfont{DejaVu Sans}}

\usepackage[french]{babel}
\usepackage[margin=1.4cm]{geometry}
\usepackage{amsmath}
\usepackage{setspace}
\usepackage{enumitem}
\usepackage{array}

% Interligne 1,5
\onehalfspacing

% Espacement réduit autour des égalités successives
\setlength{\abovedisplayskip}{1pt}
\setlength{\belowdisplayskip}{1pt}
\setlength{\abovedisplayshortskip}{0pt}
\setlength{\belowdisplayshortskip}{0pt}

% ----- Pas de numérotation des pages -----
\pagestyle{empty}

\setlength{\parindent}{0pt}

% ----- Commandes personnalisées -----
\newcounter{exercice}
\newcommand{\exercice}{%
  \refstepcounter{exercice}%
  \vspace{0.4em}
  \noindent\underline{\textbf{Exercice \arabic{exercice}}}\par
  \vspace{0em}
}

\setlist[enumerate,1]{label=\alph*),itemsep=0pt,topsep=1pt}
\setlist[itemize,1]{leftmargin=1.6em,itemsep=0pt,topsep=1pt}

\begin{document}

% ----- En-tête -----
\begin{center}
  \underline{\Large\textbf{Correction du contrôle séquence 13 version 3}}\\[0.3em]
 
\end{center}

% =====================================================================
\vspace{3cm}\exercice
\begin{enumerate}
  \item \textbf{Oui.} $56  = 8 \times 7$ 
  \newline 8 est donc un diviseur de 56.
  \item \textbf{Non.} $200 \div 9 \approx 22{,}22$
  \newline le quotient n'est pas un entier,
        donc 200 n'est pas un multiple de 9.
  \item \textbf{Non.}  $14 < 42$, donc 14 n'est pas un multiple de 42.
  \item \textbf{Non.} $50 \div 9 \approx 5{,}56$
  \newline le quotient n'est pas un entier,
        donc 9 n'est pas un diviseur de 50.
\end{enumerate}

% =====================================================================
\vspace{1cm}\exercice
\begin{enumerate}
  \item \textbf{2 et 3} (acceptés aussi : 5 ; 7 ; 11 ; \dots).
        Un nombre premier a exactement deux diviseurs : 1 et lui-même.
  \item \textbf{4 et 6} (acceptés aussi : 1 ; 8 ; 9 ; \dots).
      
  \item Dans 10 ; 21 ; 5 ; 33 ; 55 ; 68 ; 89 ; 87, les deux nombres premiers sont
        \textbf{5} et \textbf{89}.
       
\end{enumerate}

% =====================================================================
\vspace{1cm}\exercice
\textbf{a) } $10 = 2 \times 5$, 
donc les diviseurs de 10 sont $\boxed{1\ ;\ 10\ ;\ 2\ ;\ 5\ }$.

\textbf{b) } $50 = 2 \times 5 \times 5$,
donc les diviseurs de 50 sont
$\boxed{1\ ;\ 50\ ;\ 2\ ;\ 5\ ;\ 10\ ;\ 25\ }$.

% =====================================================================
\vspace{1cm}\exercice


\noindent\begin{minipage}[t]{0.48\textwidth}
\textbf{a)}
\begin{align*}
22 &= 2 \times 11
\end{align*}

\textbf{b)}
\begin{align*}
70 &= 2 \times 35\\
   &= 2 \times 5 \times 7
\end{align*}
\end{minipage}\hfill
\begin{minipage}[t]{0.48\textwidth}
\textbf{c)}
\begin{align*}
132 &= 2 \times 66\\
    &= 2 \times 2 \times 33\\
    &= 2 \times 2 \times 3 \times 11\\
\end{align*}

\textbf{d)}
\begin{align*}
195 &= 3 \times 65\\
    &= 3 \times 5 \times 13
\end{align*}
\end{minipage}

% =====================================================================
\exercice
\noindent\begin{minipage}[t]{0.31\textwidth}
\textbf{a)} 
\begin{align*}
\dfrac{4}{20} &= \dfrac{4 \times 1}{4 \times 5}\\
              &= \dfrac{1}{5}
\end{align*}
\end{minipage}\hfill
\begin{minipage}[t]{0.31\textwidth}
\textbf{b)} 
\begin{align*}
\dfrac{12}{18} &= \dfrac{6 \times 2}{6 \times 3}\\
               &= \dfrac{2}{3}
\end{align*}
\end{minipage}\hfill
\begin{minipage}[t]{0.31\textwidth}
\textbf{c)} 
\begin{align*}
\dfrac{75}{105} &= \dfrac{15 \times 5}{15 \times 7}\\
                &= \dfrac{5}{7}
\end{align*}
\end{minipage}

% =====================================================================
\vspace{1cm}\exercice
Le nombre de paquets doit diviser à la fois 90 et 105 : c'est un diviseur commun à
90 et 105. Pour en faire un maximum, on cherche le plus grand diviseur commun à 90 et 105.

Avec les décompositions en facteurs premiers :
$90 = 2 \times 3 \times 3 \times 5$ et $105 = 3 \times 5 \times 7$.
Le plus grand diviseur commun à 90 et 105 est le produit des facteurs premiers communs :
3 \times 5
                      = 15

Hulk et Batman peuvent donc faire au maximum \textbf{15 paquets},
chacun contenant  6 cookies et 7 tartelettes.

% =====================================================================
\vspace{1cm}\exercice
On utilise les critères de divisibilité :
\newline \textbf{par 2} : le chiffre des unités est 0, 2, 4, 6 ou 8 ;
\newline \textbf{par 3} : la somme des chiffres est divisible par 3
(sommes des chiffres : $23$ pour 483\,260 ; $27$ pour 333\,333\,333 ; $29$ pour 367\,445) ;
\newline \textbf{par 5} : le chiffre des unités est 0 ou 5 ;
\newline \textbf{par 10} : le chiffre des unités est 0.

\renewcommand{\arraystretch}{1.4}
\setlength{\tabcolsep}{2pt}
\begin{center}
\small
\begin{tabular}{|c|c|c|c|c|}
\hline
 & \textbf{Divisible par 2} & \textbf{Divisible par 3} & \textbf{Divisible par 5} & \textbf{Divisible par 10} \\
\hline
483\,260 & Oui & Non & Oui & Oui \\
\hline
333\,333\,333 & Non & Oui & Non & Non \\
\hline
367\,445 & Non & Non & Oui & Non \\
\hline
\end{tabular}
\end{center}

\end{document}
